% ----------------------------------------------------------
\chapter{Introdução}
% ----------------------------------------------------------

	Durante o fim do inverno austral em 2024 o estado de Santa Catarina, Brasil registrou fenômenos atmosféricos sem precedentes associados à qualidade do ar e poluição, tais como a queda na umidade relativa que em conjunto com altas concentrações de fumaça no ar criaram condições mais favoráveis para o agravamento de doenças respiratórias (as quais já estavam em seu pico por conta da estação do ano), conforme indicado pela \autocite{santacatarina_ima2026}.  As notas emitidas fizeram referência às quedas na qualidade do ar observadas nas quatro estações meteorológicas instaladas nas regiões central e litorânea do estado, enquanto as previsões meteorológicas oficiais indicavam uma atmosfera estável pelos próximos dias, fator que contribui fortemente para a estagnação de poluentes na atmosfera. Adicionalmente, registros de chuva preta nas cidades do estado mostram que o impacto não se limitou às áreas onde observações in situ puderam ser feitas  \autocite{epagri2024chuvapreta}.
	
	Neste contexto, pesquisas foram conduzidas com o intuito de avaliar o transporte de aerossóis em municípios de estados do Rio Grande do Sul e de São Paulo durante essa mesma época do ano, onde os autores apontaram os incêndios generalizados na floresta amazônica que estavam ocorrendo na época como fonte dos aerossóis transportados até as  localidades analisadas \autocite{forgioni2025}. No referente à Santa Catarina, a ausência de instrumentos de observação na região oeste do estado na época do ocorrido cria uma lacuna de informações sobre o real impacto desse episódio na população. Portanto, é proposto no presente trabalho um estudo de caso sobre o transporte e a dispersão dos aerossóis presentes na atmosfera da região oeste do estado de Santa Catarina a partir de registros feitos tanto com satélites quanto com dados de reanálise, que combinam dados observados com interpolações de saídas de modelos numéricos para cobrir áreas sem registro. Uma vez obtidos e tratados, esses dados servem de base para simulações computacionais que forneçam detalhes precisos sobre o transporte e a dispersão dos poluentes ao longo do tempo. A seguir estão dispostos tanto o objetivo geral quanto os objetivos específicos do trabalho. 



% ----------------------------------------------------------
\section{Objetivos}
% ----------------------------------------------------------

Nas seções abaixo estão descritos o objetivo geral e os objetivos específicos deste TCC.

% ----------------------------------------------------------
\subsection{Objetivo Geral}
% ----------------------------------------------------------

	Simular o transporte e dispersão dos poluentes presentes na região oeste de Santa Catarina durante os dias de máxima profundidade óptica de aerossóis observadas em setembro de 2024.

% ----------------------------------------------------------
\subsection{Objetivos Específicos}
% ----------------------------------------------------------

	\begin{itemize}
		\item Identificar o padrão atmosférico da época sobre a América do Sul, com ênfase nos sistemas e mecanismos que influenciaram a intensidade e duração do evento.
		\item Analisar o tipo e a concentração de aerossóis que foram registrados no estado de Santa Catarina.
		\item Gerar simulações de trajetória e dispersão no Modelo Hysplit a partir de dados de reanálise de variáveis meteorológicas e de emissão por queimadas.
	\end{itemize}
	
	
	